\subsection{Pairing of parentheses}
\begin{lem} \label{lem.pairing.1}
A proper pairing of $2n$ parentheses contains (at least) an innermost pair. More
precisely, if we have a proper pairing for $2n$ parentheses, then there exists
an inner pair in this pairing.
\end{lem}
\begin{proof}
Let $P(n)$ denote the lemma. Let $Q(m)$ denotes that $P(n)$ holds for all $n
\le m$. To apply course-of-values induction, we prove first $P(0)$ and then
$Q(m) \to Q(m+1)$.

$P(0)$ is vacuously true.

Suppose inductively $Q(m)$ holds. Let's consider a pair of $2(m+1)$
parentheses. If it can admit a proper pairing at all, then the leftmost
parenthesis must be a left one to not violate the rule: ``a left parenthesis is
always paired with ... to the right of it.'' Although the same reasoning
applies to deduce that the rightmost must be a right parenthesis, it needs not
to match the leftmost one. 

In any case, the leftmost one much be matched by some parathesis to the right
of it. If they are immediate next to each other, then we have found such an
innermost pair. Otherwise, some other paratheses are inside. Because a proper
pairing is admitted by the whole, those inside must also satisfy the rule of
proper pairing and admit a smaller one themselves. Moreover, they are of $\le
2m$ parentheses. Thus, appkying the inductive hypothesis, we know they have an
innermost pair, as desired.
\end{proof}

\begin{lem}
A set of $2n$ parentheses admits at most one proper pairing. In other words,
two proper pairings of $2n$ parentheses must be the same.
\end{lem}
\begin{proof}
Induct on $n$. When $n=0$, it's vacuously true.

Let $p_1$ and $p_2$ be two proper pairings on $2n$ parentheses. By
Lemma~\ref{lem.pairing.1}, we must see an innermost pair $i_1$ for $p_1$ and
$i_2$ for $p_2$. 

There are two cases: If we can choose $i_1$ and $i_2$ to be the same, then we
choose so, and exclude it from both $p_1$ and $p_2$. Then,
the whole with $i_1$ or $i_2$ removed must be the same, which can only admit at
most one proper pairing, by the inductive hypothesis. The only possible proper
pairing of the whole is the outer plus the innermost as a pair, for we cannot
make a proper pairing if any one of the innermost pair is assigned to one of
the outer. Thus, $p_1$ and $p_2$ are the same.

If we cannot choose $i_1$ and $i_2$ to be the same, then $i_1$ is not assigned
as a pair in $p_2$ and vice versa. However, such an assignment cannot be part
of a proper pairing, as then either $i_1$ or $i_2$ will be part of separated
pairs.

We can now close the induction.
\end{proof}

\begin{lem}
If $2n$ parentheses and a consecutive subset of $2m$ of them both admit a
proper pairing, then they must agree with each other. 
\end{lem}
\begin{proof}
We induct on $m$. When $m=0$, it's vacuously true.

Inductively suppose that it's already true for $m$. Consider a subset of $2m+2$
parentheses. By Lemma~\ref{lem.pairing.1}, we can remove the innermost pair of
it. Moreover, in both the $2m$'s and the $2n$'s proper pairing, the inner pair
must pair itself. And now we have the $2m$'s as a subset, which by the
inductive hypothesis agrees with the $2n$'s. 

Adding the removed innermost pair back, we can now close the induction.
\end{proof}
